\section{Ajuste y convergencia del modelo}

\subsection{Resultados globales de estimación}

La Tabla~\ref{tab:ajuste_convergencia} resume las métricas globales del modelo principal
territorial bajo las dos formas de parametrizar las variables comunes. La Opción A mantiene a
\texttt{BIP} como alternativa de referencia para las variables comunes, mientras que la Opción B
usa una restricción de suma cero para reportar coeficientes para las tres alternativas. Ambas
especificaciones representan el mismo contenido conductual y, por lo tanto, entregan el mismo
ajuste global.

\begin{table}[htbp]
\centering
\small
\caption{Ajuste y convergencia del modelo principal territorial}
\label{tab:ajuste_convergencia}
\begin{tabular}{lrr}
\toprule
Métrica & Opción A & Opción B \\
\midrule
Parámetros estimados & 52 & 52 \\
Observaciones & 466{,}639 & 466{,}639 \\
Observaciones excluidas & 0 & 0 \\
Log-verosimilitud inicial & -512{,}655.3 & -512{,}655.3 \\
Log-verosimilitud final & -232{,}257.1 & -232{,}257.1 \\
Test de razón de verosimilitud & 560{,}796.5 & 560{,}796.5 \\
Rho-cuadrado & 0.5470 & 0.5470 \\
Rho-cuadrado ajustado & 0.5470 & 0.5470 \\
AIC & 464{,}618.2 & 464{,}618.2 \\
BIC & 465{,}193.0 & 465{,}193.0 \\
Norma final del gradiente & 4.7659 & 4.8173 \\
Convergencia & Sí & Sí \\
\bottomrule
\end{tabular}
\end{table}

En ambos casos, el modelo converge correctamente y no excluye observaciones de la muestra de
estimación. La igualdad en log-verosimilitud, AIC y BIC confirma que las dos opciones no
corresponden a modelos sustantivamente distintos, sino a dos normalizaciones alternativas para
presentar los coeficientes de variables comunes. La diferencia principal está en la interpretación
de los parámetros: en la Opción A los coeficientes de variables comunes se leen respecto de
\texttt{BIP}, mientras que en la Opción B los coeficientes se leen como desviaciones respecto del
promedio de alternativas.

\subsection{Lectura de las métricas de ajuste}

La log-verosimilitud final mejora sustancialmente respecto de la log-verosimilitud inicial, pasando
de $-512{,}655.3$ a $-232{,}257.1$. Esto indica que las variables incluidas en la especificación
aportan capacidad explicativa relevante respecto de un modelo sin parámetros informativos. En la
misma línea, el test de razón de verosimilitud alcanza un valor de $560{,}796.5$, lo que refleja una
mejora global importante frente al modelo inicial.

El valor de rho-cuadrado es $0.5470$, lo que sugiere un ajuste alto para un modelo de elección
discreta aplicado a viajes individuales. Esta métrica no debe interpretarse como un $R^2$ de
regresión lineal, pero sí permite evaluar cuánto mejora el modelo respecto de una especificación
base. En este caso, el resultado indica que el modelo logra capturar una parte sustantiva de la
estructura observada de elección entre \texttt{BIP}, \texttt{QR\_RED} y \texttt{QR\_OTHER}.

Los criterios AIC y BIC son iguales entre ambas opciones, porque ambas tienen la misma
log-verosimilitud y el mismo número de parámetros estimados. Por lo tanto, no sirven para escoger
entre la Opción A y la Opción B. La elección entre ambas debe entenderse como una decisión
metodológica de interpretación y presentación de resultados, no como una comparación de ajuste
predictivo o estadístico.

Finalmente, ambas especificaciones reportan convergencia numérica. La Opción B requiere más
iteraciones que la Opción A debido a la parametrización con restricción de suma cero, pero alcanza
el mismo óptimo de log-verosimilitud. Por esta razón, la decisión entre ambas rutas se deja
planteada como una decisión de reporte: mantener una lectura tradicional respecto de \texttt{BIP} o
presentar coeficientes comparables para las tres alternativas.
